%!TEX encoding = UTF8
%!TEX root = 0-notes.tex

\chapter{Espaces vectoriels}
\label{chap:espaces-vectoriels}

Produit matriciel $Ax$ du point de vue vectoriel $Ax = x_i a_i$.

\section{Définition et exemples}

% ev, vecteur
% combinaison linéaire
% colinéarité

% R[x]
% droites, plans
% positions relatives
% intersections

\section{Indépendance, base}

\subsection{Base orthonormale}

\dfn{indépendance linéaire}{
	
}{dfn:indep-lin}

\dfn{base}{

}{dfn:base}

\nomen{
	Une base dont les vecteurs sont orthogonaux deux à deux et dont chaque vecteur est de norme 1 est une \emphindex{base orthonormée}.
}

\ex{base canonique}{
	
}{ex:base-canonique}

\nt{
	Il s'avère que toute base de $\R^2$ est de taille 2 : c'est sa \emphindex{dimension} en tant qu'espace vectoriel sur $\R$.
	
	Plus généralement, la dimension de $\R^n$ comme \emphindex{$\R$-espace vectoriel} est égale à $n$.
}

% indépendance
% base

\section{Espaces affines}

\section*{Exercices supplémentaires}
\addcontentsline{toc}{section}{Exercices supplémentaires}

\shipoutExercise

\section*{Solutions}
\addcontentsline{toc}{section}{Solutions}

\shipoutAnswer



